\jlreqsetup{warichu_opening=, warichu_closing=}

\usepackage{ifthen}

\usepackage{mdframed}

\usepackage{graphicx}

\usepackage{lltjp-geometry}

\usepackage[driver=dvipdfmx]{geometry}

\usepackage{pxrubrica}

\usepackage{ulem}

\usepackage{luatexja-fontspec}

\usepackage{newunicodechar}

\usepackage{caption}

\usepackage[kumi=aki, tateaki=1]{kanbun}

%%%%%%%%%%%%%%%%%%%%%%%%%%%%%%
% environments
%%%%%%%%%%%%%%%%%%%%%%%%%%%%%%

% \begin{azconvIndented}{n}{m}はｍ字下げ、改行してｎ字下げ。
% ［＃ｎ字下げ］はｎ字下げ、改行してｎ字下げ。
% ［＃改行天付き］はｎ字下げ、改行して０字下げ。
\newenvironment{azconvIndented}[2]{
  \setlength{\parindent}{#2}
   \leftskip=#1

  }
  {
   \setlength{\parindent}{0\zw}
   \leftskip=0em
  }

% \begin{azconvAlignBottom}{n}は下からｎ字上げ。地付きは下から０字上げ。 
\newenvironment{azconvAlignBottom}[1]{
  \rightskip=#1
 }
 {\rightskip=0\zw}

% #1は字下げ及び字詰めのスコープ内の場合指定される。字下げの場合下げる字数。
% 字詰めの場合は求められる字数に近づくように設定される。
 \newenvironment{azconvFramed}[1][]%
 {\begin{mdframed}[#1,skipabove=8pt]}
{\end{mdframed}}


\newenvironment{azconvNarrow}[1]{
 \rightskip= \dimexpr\linewidth-#1}
{\rightskip=0em}

\newenvironment{azconvFigure}{
 \begin{figure}[h!]
\begin{center}}
{\end{center}
\end{figure}}


%%%%%%%%%%%%%%%%%%%%%%%%%%%%%%
% commands
%%%%%%%%%%%%%%%%%%%%%%%%%%%%%%

%%%%% (re)define various under/overline styles for ulem
%%%%% shamelessly reusing and adjusting code from ulem.sty

\makeatletter

\UL@protected\def\uline{\leavevmode \bgroup
 \UL@setULdepth
 \ifx\UL@on\UL@onin \advance\ULdepth2.8\p@\fi
 \advance\ULdepth-0.5\p@\markoverwith{\lower\ULdepth\hbox
   {\kern-.03em\vbox{\hrule width.2em}\kern-.03em}}%
 \ULon}

\UL@protected\def\oline{\leavevmode \bgroup
 \UL@setULdepth
 \ifx\UL@on\UL@onin \advance\ULdepth2.8\p@\fi
 \advance\ULdepth-0.5\p@\markoverwith{\raise\ULdepth\hbox
   {\kern-.03em\vbox{\hrule width.2em}\kern-.03em}}%
 \ULon}

 \UL@protected\def\uuline{\leavevmode \bgroup
 \UL@setULdepth
 \ifx\UL@on\UL@onin \advance\ULdepthi2.8\p@\fi
 \advance\ULdepth1\p@\markoverwith{\lower\ULdepth\hbox
   {\kern-.03em\vbox{\hrule width.2em\kern1\p@\hrule}\kern-.03em}}%
 \ULon}

\UL@protected\def\ooline{\leavevmode \bgroup
 \UL@setULdepth
 \ifx\UL@on\UL@onin \advance\ULdepth2.8\p@\fi
 \advance\ULdepth-0.5\p@\markoverwith{\raise\ULdepth\hbox
   {\kern-.03em\vbox{\hrule width.2em\kern1\p@\hrule}\kern-.03em}}%
 \ULon}

\UL@protected\def\ooline{\leavevmode \bgroup
 \UL@setULdepth
 \ifx\UL@on\UL@onin \advance\ULdepth2.8\p@\fi
 \advance\ULdepth-.8\p@\markoverwith{\raise\ULdepth\hbox
   {\kern-.03em\vbox{\hrule width.2em\kern1\p@\hrule}\kern-.03em}}%
 \ULon}

\UL@protected\def\dotuline{\leavevmode \bgroup 
  \UL@setULdepth
  \ifx\UL@on\UL@onin \advance\ULdepth2.8\p@\fi
  \markoverwith{\begingroup
   \advance\ULdepth-2\p@\lower\ULdepth\hbox{\kern.06em .\kern.04em}%
     \endgroup}%
  \ULon}

\UL@protected\def\dotoline{\leavevmode \bgroup \UL@setULdepth
 \ifx\UL@on\UL@onin \advance\ULdepth2.8\p@\fi
 \markoverwith{\begingroup
 \advance\ULdepth3\p@\raise\ULdepth\hbox{\kern.06em .\kern.04em}%
	\endgroup}%
 \ULon}

\UL@protected\def\dashuline{\leavevmode \bgroup 
  \UL@setULdepth
  \ifx\UL@on\UL@onin \advance\ULdepth2\p@\fi
  \markoverwith{\begingroup
     \advance\ULdepth-0.8\p@\ 
     \lower\ULdepth\hbox{\kern.06em -\kern.04em}%
     \endgroup}%
  \ULon}

\UL@protected\def\dasholine{\leavevmode \bgroup \UL@setULdepth
 \ifx\UL@on\UL@onin \advance\ULdepth2\p@\fi
 \markoverwith{\begingroup
	\advance\ULdepth1.5\p@\raise\ULdepth\hbox{\kern.06em -\kern.04em}%
	\endgroup}%
  \ULon}

\UL@protected\def\owave{\leavevmode \bgroup 
  \ifdim \ULdepth=\maxdimen \ULdepth 3.5\p@
  \else \advance\ULdepth2\p@ 
  \fi \advance\ULdepth3.2\p@\markoverwith{\raise\ULdepth\hbox{\sixly \char58}}\ULon}
\font\sixly=lasy6 % does not re-load if already loaded, so no memory drain.

\UL@protected\def\uwave{\leavevmode \bgroup 
  \ifdim \ULdepth=\maxdimen \ULdepth 3.5\p@
  \else \advance\ULdepth2\p@ 
  \fi \advance\ULdepth0.5\p@\markoverwith{\lower\ULdepth\hbox{\sixly \char58}}\ULon}
\font\sixly=lasy6 % does not re-load if already loaded, so no memory drain.

\makeatother

%********


\newcommand{\azconvBold}[1]{%
{\textbf{#1}}}

\newcommand{\azconvContributor}[1]{
   \hfill #1

  }

\newcommand{\azconvFbox}[1]{
   \fbox{#1}
  }

\newcommand{\azconvIncludegraphics}[3]{
 \begin{center}
  \hskip -4\zw
  \resizebox{#2}{#1}{\includegraphics[angle=90, origin=c]{#3}}
\end{center}}

\newcommand{\azconvIncludegraphicsInline}[1]{
 \raisebox{-0.6\zw}{\includegraphics[angle=90, origin=c, width=1\zw, height=1\zw]
{#1}}}

\newcommand{\azconvInlineHeader}[2]{%
\hskip 0.5\zw\textbf{#2}\hskip-0.5\zw}

\newcommand{\azconvItalic}[1]{%
{\textit{#1}}}

\newcommand{\azconvKenten}[3][]{%
 \ifthenelse{\equal{#1}{left}}%
 {\kentenmarkintate{#2}\kenten[S]{#3}}%
 {\kentenmarkintate{#2}\kenten{#3}}%
}

\newcommand{\azconvKunten}[1]{%
\raisebox{-0.4\zw}{\scriptsize #1}\normalsize}

%　#1は改ページ、改丁など注記で指示されているもの。デフォルトでは無視する。
\newcommand{\azconvNewpage}[1]{%
 \newpage
}

\newcommand{\azconvOkurigana}[1]{%
\raisebox{0.6\zw}{\scriptsize #1}\normalsize}

\newcommand{\azconvSection}[1]{%
 \section{#1}
}

\newcommand{\azconvSubsection}[1]{%
  \subsection{#1}
 }

\newcommand{\azconvSubsubsection}[1]{%
  \subsubsection{#1}
 }


\newcommand{\azconvSup}[1]{%
\textsuperscript{#1}}

\newcommand{\azconvSub}[1]{%
\textsubscript{#1}}

\newcommand{\azconvSubtitle}[1]{
	{\large
		\begin{azconvIndented}{7\zw}{0\zw}

			#1

	 \end{azconvIndented}
	}
 
  }

\newcommand{\azconvTatechuyoko}[1]{%
 \tatechuyoko{#1}}

\newcommand{\azconvTitle}[1]{
	{\LARGE
	 
	 \begin{azconvIndented}{3\zw}{0\zw}

	 #1
	 
	\end{azconvIndented}
		\vskip 1\zw
	}

  }

 % #1が傍線の種類、#2が左右、#3が対象文字列。
\newcommand{\azconvUline}[3]{%
   \ifthenelse
   {\equal{#1}{solid}}
	{\ifthenelse{
	 \equal{#2}{under}}
	{\oline{#3}}
   {\uline{#3}}}{}%
   %
\ifthenelse
   {\equal{#1}{double}}
	{\ifthenelse{
	 \equal{#2}{under}}
	 {\ooline{#3}}
	{\uuline{#3}}}{}%
	%
\ifthenelse
   {\equal{#1}{dotted}}
	{\ifthenelse{
	 \equal{#2}{under}}
	 {\dotoline{#3}}
	{\dotuline{#3}}}{}%
	%
\ifthenelse
   {\equal{#1}{dashed}}
	{\ifthenelse{
	 \equal{#2}{under}}
	 {\dasholine{#3}}
	{\dashuline{#3}}}{}%
	%
\ifthenelse
   {\equal{#1}{wavy}}
	{\ifthenelse{
	 \equal{#2}{under}}
	 {\owave{#3}}
	{\uwave{#3}}}{}%
}

% #1は必要となるboxのサイズ、#2が文字列。#1はjlreqの場合自動で算出されるので無視。
\newcommand{\azconvWarichu}[2]{
 \warichu{#2}}

% 割中の中で改行が指定されている場合に使う。#1はboxのサイズ、#2が文字列。
\newcommand{\azconvWarichuM}[2]{\linespread{0.5}
{\scriptsize\parbox{#1}{#2}}}

\newcommand{\azconvYokogumi}[1]{%
\raisebox{-.35\zw}{\hbox{\utod #1}}}


 %%%%%%%%%%%%%%%%%%%%%%%%%%%%%%
% for Kanbun kunten 
%
% taken from 漢文マクロ kambunsnx 
%
% 峰山進(閑舎) https://rakunet.org/p/2022/0729
%
%%%%%%%%%%%%%%%%%%%%%%%%%%%%%%

\def\azconvIR{一\kern-.6\zw レ}
\def\azconvNR{二\kern-.4\zw レ}
\def\azconvJR{上\kern-.25\zw レ}
\def\azconvKR{甲\kern-.4\zw レ}


%%%%%%%%%%%%%%%%%%%%%%%%%%%%%%
% setup basic lengths, fonts, etc.
%%%%%%%%%%%%%%%%%%%%%%%%%%%%%%

\setlength{\parindent}{0em}

\renewcommand{\baselinestretch}{1.2}

\captionsetup{font={small, normalfont},labelformat=empty,skip=0pt}

\ModifyHeading{section}{indent=2\zw, format={\Large #2}}
\ModifyHeading{subsection}{indent=3\zw,format={\large #2}}
\ModifyHeading{subsubsection}{indent=4\zw,format={\bfseries #2}}

%\usepackage[noto-jp]{luatexja-preset}

%%%%%%%%%%%%%%%%%%%%%%%%%%%%
%
% Fallback for unavailable characters
%
%%%%%%%%%%%%%%%%%%%%%%%%%%%%


%\newjfontfamily{\fallback}{MingLiU-ExtB}

%\newunicodechar{𥥔}{{\fallback 𥥔}} 
